\section{Methodology}

At its core, adapting an English word to Filipino orthography is an instance of
\emph{orthographic nativization}: the process by which a loanword is respelled
to conform to the recipient language's writing conventions.
Filipino orthography is fundamentally phonemic---it follows the principle
\emph{kung ano ang tunog, yun ang sulat} (what you hear is what you write),
where each letter maps to a single sound with near-perfect
correspondence. English orthography, by contrast, is
\emph{deep}: a single letter may represent multiple phonemes, and a single
phoneme may be spelled in multiple ways. The task, then, is to bridge a deep
orthography and a shallow one---to transform a sequence of English
graphemic units into a sequence of Filipino graphemic units that faithfully
represents how the word is perceived and pronounced in Filipino.

\subsection{Graphemic Tokenization}

Before any mapping can take place, the English input must be segmented into
its functional graphemic units. This step is necessary because English
orthography contains multi-letter patterns that function as single phonological units:
\begin{itemize}
    \item \textbf{Digraphs}: \emph{ph}, \emph{ch}, \emph{sh}, \emph{th}, \emph{ng}
    \item \textbf{Trigraphs}: \emph{igh}, \emph{tch}, \emph{dge}
    \item \textbf{Higher-order multigraphs}: \emph{ough}, \emph{eigh}, \emph{augh}
\end{itemize}
For instance, \emph{ch} in the word ``chapter'' represents a single phonemic
unit, the voiceless postalveolar affricate /\textipa{tS}/, and must be treated
as one token rather than two independent letters.

The tokenizer applies a greedy, longest-match strategy: at each position in the
input string, it selects the longest graphemic pattern recognized by the system.
For example, the words ``chapter'', ``straight'', and ``phishing'' are tokenized
as:
\begin{center}
    \texttt{Ch~A~P~T~Er} \\
    \texttt{S~T~R~Aigh~T} \\
    \texttt{Ph~I~Sh~I~Ng}
\end{center}
This tokenization ensures that the subsequent mapping rules operate on
phonologically meaningful units rather than on raw characters.

\subsection{Rule Structure}

The mapping from English grapheme tokens to their Filipino orthographic
equivalents is governed by a set of rules. In constructing these rules, we
observe that English graphemes vary considerably in the difficulty they pose
for adaptation.

Most English consonant graphemes have a single or highly predictable
pronunciation. Their adaptation to Filipino is straightforward and follows
well-established nativization patterns. For example, the grapheme \emph{b}
consistently represents the voiced bilabial plosive /\textipa{b}/ and maps
directly to Filipino \emph{b}. Graphemes for sounds absent from native Filipino
phonology follow regular substitution conventions: \emph{v}~$\rightarrow$~\emph{b}
(e.g., ``volleyball'' $\rightarrow$ ``balibol''), \emph{f}~$\rightarrow$~\emph{p}
(e.g., ``traffic'' $\rightarrow$ ``trapik''), and
\emph{th}~$\rightarrow$~\emph{t} (e.g., ``three'' $\rightarrow$ ``tri'').
These are consistent with the phonemic inventory of Filipino, which historically
lacked /f/, /v/, and /\textipa{T}/.

Vowels, however, are far more problematic. English is rich in vowel contrasts,
and its orthography encodes these distinctions poorly. The grapheme \emph{a}
alone may represent /\textipa{\ae}/ in ``cat,'' /\textipa{eI}/ in ``cake,''
/\textipa{A:}/ in ``father,'' and /\textipa{@}/ in ``about''. Filipino, with
only five vowels (/a/, /e/, /i/, /o/, /u/), must collapse these
distinctions. Furthermore, because English is a stress-timed language, unstressed
vowels are frequently reduced to the schwa /\textipa{@}/, adding another layer
of ambiguity.

These observations motivate a three-tiered rule system:

\subsubsection{Context-Free Orthographic Rules}

Context-free rules map a grapheme token to its Filipino equivalent in
isolation, without reference to surrounding tokens. Formally, we can express
such a rule as a function
$f(\omega_1, \omega_2, \dots, \omega_n,\; i) \mapsto \phi_i$
that depends only on the identity of $\omega_i$, not on its neighbors.
Consonant graphemes with a unique and predictable pronunciation belong to this
category. For instance, the grapheme \emph{b} always maps to \emph{b};
hence, ``basketball'' becomes ``basketbol''.

\subsubsection{Context-Sensitive Orthographic Rules}

Context-sensitive rules consider the graphemes surrounding the target token.
These are necessary when a grapheme's pronunciation---and therefore its Filipino
adaptation---depends on its orthographic environment.

Consider the grapheme \emph{p} in the word ``psychology''. In isolation,
\emph{p} would map to Filipino \emph{p}; however, when followed by \emph{s} at
the word onset, it is silent. Examining the neighboring graphemes allows the
rule to correctly omit \emph{p}, yielding ``saykoloji'' rather than ``psaykoloji''. Similarly, the grapheme \emph{c} maps to /\textipa{s}/ before
front vowels (\emph{e}, \emph{i}, or \emph{y}) but to /\textipa{k}/ elsewhere, requiring
context to resolve.

\subsubsection{Phonetic Rules}

As discussed, vowels resist prediction from orthographic context alone. Because
Filipino orthography requires that each vowel letter correspond to a definite
vowel sound, we cannot avoid determining the actual phonetic quality of each
vowel. The algorithm therefore incorporates phonetic information---such as
lexical stress, vowel quality, and reduction patterns---to resolve cases where
orthographic rules are insufficient. Stress placement is especially critical:
stressed vowels retain their full quality and must be mapped to the appropriate
Filipino vowel (/a/, /e/, /i/, /o/, or /u/), while unstressed vowels reduced
to schwa must also be assigned a Filipino vowel, typically on the basis of the
English spelling or the closest perceptual match.